\cotefile{157-5}
\scanpage{373}
\mspage{289}
% archive sheet number: 372

\uline{4.13.5.} Soit $\underline S$ une sous-catégorie \uuline{str. pleine}
de $M$ (catégorie accessible), \inserted{$\pi'_0$ un cardinal infini.}
Dire que $\underline S$ est $\pi'_0$-accessible dans $M$ signifie que
les conditions suivantes sont satisfaites

1) $\pi'_0$ est utile pour $\underline S$, i.e. $\underline S \overset{\sim}{\underset{\sim}{\leftarrow}}\operatorname{Ind}_{\pi'_0}(S_{\pi'_0})$\unsure{;}

2) $\underline S\hookrightarrow M$ est $\pi'_0$-accessible

3) $\underline S_{\pi'_0}\subset M_{\pi'_0}$ (qui implique que $S_{\pi'_0}$ est
équivalente à une petite catégorie).

\deleted{Soit} Supposons ces conditions satisfaites,
et soit $\pi^{\unsure{\prime|\diamond}}$ un cardinal $>\pi'_0$. Est-ce que
$\underline S$ est $\pi^{\unsure{\prime}}$-accessible \inserted{dans $M$}, i.e. les conditions
similaires sont-elles satisfaites pour $\pi$?
C'est évident pour $2^\circ$), mais risque d'être
faux pour $1^\circ$, $3^\circ$.

\deleted{\deleted{C'est pourquoi il a été utile \unsure{d'introduire|de déterminer}} le plus petit cardinal
$\pi'_0$ \inserted{\unsure{infini}} tel que $\pi'_0$ soit utile pour $\underline S$, et
que $\underline S\hookrightarrow M$ soit $\pi'_0$-accessible, et
$\pi'_1$ \deleted{$=\operatorname{Sup}(\pi'_0,\operatorname{card}\unsure{\operatorname{tot}|\operatorname{Fl}}S_{\pi'_0})$} le plus petit
cardinal $\geq$\inserted{$\pi'_0$,} $(\pi'_0,\operatorname{card}\operatorname{tot}S_{\pi'_0}$, et tel que
$S_{\pi'_0}\subset\deleted{\unsure{\cong}}M_{\pi'_1}$. \deleted{Si alors} (NB $\pi_0$ est le plus
petit cardinal \deleted{tel que} utile pour $M$.)}
\deleted{Si alors $\pi\geq\pi'_1$} Mais faisons appel à
% slanted
\margin{\deleted{\unsure{En fait, distinguons} \unsure{rel.} \unsure{à} \illegible{1} \unsure{4.13.1} à \unsure{4.13.2} \unsure{ici}.} \unsure{4} \unsure{4.}}
4.13.2. \unsure{c)}. \deleted{Supposons que $\pi'_0\geq\pi_0$,} Posons
\[
\pi'_1=\operatorname{Sup}(\pi'_0,\operatorname{card}\operatorname{tot} S_{\pi'_0},\pi_0)
\]
\margin{\unsure{$\Vert$|``}plus petit cardinal infini \deleted{\unsure{tel que}} utile pour $M$\unsure{.}}
\scanpage{374}
\mspage{290}
% archive sheet number: 373

Alors \unsure{pour} $\pi\geq\pi'_1$ tel que $\pi^{\pi'_0}=\pi$,
$\underline S$ est $\pi$-accessible dans $M$ (4.13.2. b et c).

Supposons p. ex.
\[
\pi'_0\geq\pi_0
\]
(pratiquement ce sera tjrs vérifié) et de plus
\[
\pi'_0\geq\deleted{\unsure{\pi_1=\operatorname{Sup}}}\operatorname{Card}\operatorname{tot}M_{\pi_0}
\]
i.e. supposons
\[
\pi'_0\geq\operatorname{Sup}(\pi_0,\operatorname{Card}\operatorname{tot} M_{\pi_0})\equiv\pi_1.
\]
Alors \deleted{\unsure{comme} $\unsure{\in S_{\pi'_0}}$}
\[
\deleted{\pi'_1=\operatorname{Sup}(\pi'_0,\operatorname{Card}\unsure{\operatorname{Fct}|\operatorname{tot}|\operatorname{Fl}}S_{\pi'_0})}\ \leftarrow
\]
\[
\pi'_1=\operatorname{Sup}(\pi'_0,\operatorname{card}\operatorname{tot}S_{\pi'_0})\leq\operatorname{Sup}(\pi'_0,\operatorname{card}\operatorname{tot}M_{\pi'_0})\leq 2^{({\pi'_0}^{\pi_0})}.
\]
(4.7.7.1 p. \unsure{173|177}). Donc on trouve ceci :

\uline{Corollaire} 4.13.6. Soit $\pi'_0$ un cardinal $\geq\pi_1$,
\inserted{$\operatorname{Sup}(\pi_0,\operatorname{card}\operatorname{tot} M_{\pi_0})$}
et soit $\underline S$ une sous-catégorie $\pi'_0$-accessible
dans $M$. Alors pour \deleted{tt} cardinal $\pi$
% slanted
\margin{\uline{NB} $\pi_0$ est le plus petit cardinal utile pour $M$}
tel que $\pi\geq 2^{({\pi'_0}^{\pi_0})}$, \inserted{et} tel que l'on
ait $\pi^{\pi'_0}=\pi$ (i.e. $\pi$ de la forme
$c^{\pi'_0}$, où $c$ un cardinal $\geq 2^{({\pi'_0}^{\pi_0})}$),
$\underline S$ est $\pi$-accessible dans $M$ ! (De plus,
ces cardinaux sont réguliers pour $\underline S$ \uuline{et}
pour $M$, donc $\underline S\approx\operatorname{Ind}_{\pi}(\underline S_{\pi})$, $M\cong\operatorname{Ind}_{\pi}(M_\pi)$.)

\uline{Enchaîner vers \unsure{nos|mes|n°s} 4.13.23. \unsure{etc.}}

% slanted
\margin{Mettre\unsure{!} ici le lemme \unsure{4.16.7.11|4.16.7.1)|4.16.7.II} signalé en marge p. 292}
\scanpage{375}
\mspage{291}
% archive sheet number: 374

% slanted
\margin{le \unsure{rejeter|reporter} après 4.13.23}
\uline{Proposition} 4.13.7. Soit $M$ une catégorie accessible,
$(\unsure{\underline S|S}_\alpha)_{\alpha\in A}$ une \deleted{ensemble}\inserted{famille} de parties de $\operatorname{Ob}M$,
indexée par un petit ensemble $A$, $(\underline S_\alpha)_{\alpha\in A}$
la famille des sous-catégories accessibles dans $M$
correspondantes. Soit $S=\bigcap_{\alpha\in A}S_\alpha$, $\underline S=\bigcap_{\alpha\in A}\underline S_\alpha$, de
sorte que $\underline S$ est la sous-catégorie pleine de
$M$ définie par $S\subset\operatorname{Ob}M$. Ceci posé, si les
$S_\alpha$ sont accessibles, il en est de même
de $S$ -- ou encore, si les $\underline S_\alpha$ sont accessibles
dans $M$, il en est de $=$ de $\underline S$.

Cela résulte du résultat plus précis :

\uline{Corollaire} 4.13.8. Avec les notations \deleted{précéd. et \unsure{de}}\inserted{de 4.13.7 et 4.13.6,}
soit, pour tt $\alpha\in A$, $\pi'_\alpha$ un cardinal \inserted{$\geq\pi_1$} tel que
$\underline S_\alpha$ soit $\pi'_\alpha$-accessible.
% slanted
\margin{Cf reformulation plus \unsure{naturelle|rationnelle} et plus courante \unsure{en} pages 302, 302 bis \unsure{et|à} \unsure{ter}.}

\deleted{Soit \deleted{$\unsure{\pi'_0=\operatorname{Sup}(\operatorname{card}A,\operatorname{Sup}_\alpha 2^{({\pi'_\alpha}^{\pi_0})})}$}
\deleted{$\unsure{\pi=\operatorname{Sup}(\operatorname{card}}$}
$\pi'_0=\operatorname{Sup}(\{\operatorname{card}A,(\unsure{\pi'_\alpha})_{\alpha\in A}\})$,
et soit $\pi\geq\unsure{2^{{\pi'_0}^{\pi_0}}}$ \deleted{\unsure{infini tel que} $\unsure{\pi\geq}$}
un cardinal tel que $\pi^{\pi'_0}=\pi$. Alors $S$
est $\pi$-accessible, i.e. $\underline S$ est $\pi$-accessible dans $M$.
On aura en effet} Soit
\[
\begin{aligned}
\pi'_0&=\deleted{\unsure{\operatorname{Sup}(\operatorname{card}A,}}\operatorname{Sup}_{\alpha\in A}\pi'_\alpha \deleted{\unsure{,\operatorname{card} A})}\\
\pi''_0&=\operatorname{Sup}(\operatorname{card}A\deleted{\unsure{,\pi}},\operatorname{Sup}_\alpha 2^{{\pi'_\alpha}^{\pi_0}})
\end{aligned}
\]
% slanted
\margin{Ça a l'air un peu \unsure{abracadabrant} ! Il faut de plus supposer que $\pi$ n'est pas cardinal \unsure{et|$\infty$|$\aleph$} limite !}
et soit $\pi\geq\pi''_0$ un cardinal tel que
$\pi^{\pi'_0}=\pi$, i.e. de la forme $c^{\pi'_0}$, avec $c\geq\pi'_0$.
Alors $S$ est $\pi$-accessible (i.e. $\underline S$ est $\pi$-accessible
dans $M$). En particulier, $S$ est $2^{\pi''_0}$-accessible.
\scanpage{376}
\mspage{292}
% archive sheet number: 375

% slanted
\margin{\deleted{Il \unsure{y a} \unsure{lieu}}}
\uline{Démonstration} \deleted{En vertu de 4.13.6,} \inserted{Soit $\alpha\in A$.} Comme
$\pi'_\alpha\leq\pi'_0$ et $\pi^{\pi'_0}=\pi$, on
\deleted{\unsure{a aussi} \inserted{$\unsure{a\ \pi\geq}$} $\pi^{\pi'_\alpha}=\pi$, $(\pi_0\geq 2^{({\pi'_\alpha}^{\pi_0})})$, $S$} a fortiori
\[
\pi^{\pi'_\alpha}=\pi,\quad\text{de plus}\quad \pi\geq\pi''_0\geq 2^{({\pi_\alpha}^{\pi_0})},
\]
donc en vertu de 4.13.6, $\underline S_\alpha$ est
$\pi$-accessible \inserted{ds $M$}. D'autre part, comme $\underline S_\alpha$ est
$\pi'_\alpha$-accessible, il est stable \inserted{dans $M$} par $\varinjlim$ grandes
devant $\pi'_\alpha$, a fortiori par $\varinjlim$ gdes devant
$\pi'_0$ ($\geq\pi'_\alpha$). D'ailleurs, $\pi^{\pi'_0}=\pi$ implique
$\pi\geq 2^{\pi'_0}>\pi'_0$. Donc \deleted{le corollaire}\inserted{4.13.8} résulte du

\uline{Corollaire} 4.13.9. Soit $\pi\unsure{\geq|>}\pi_0$ un cardinal
tel que $\underline S_\alpha$ soit $\pi$-accessible \inserted{dans $M$} pour \deleted{\unsure{tt $A$}}
$\alpha\in A$, et tel que $\operatorname{card} A\leq\pi$. Supposons
de plus qu'il \deleted{$\underline S_\alpha$ soit} existe un cardinal
$\pi'_0$ \deleted{$\unsure{\geq}$} tel que $\pi_0\leq\pi'_0<\pi$, tel que les
$\underline S_\alpha$ soient stables \unsure{dans} $M$ par $\varinjlim$ gdes
devant $\pi'_0$. Alors $\underline S=\bigcap_{\alpha\in A}\underline S_\alpha$ est $\pi$-accessible dans $M$.

% slanted
\margin{Ce serait le moment de formuler ici le lemme \unsure{4.16.7.11|4.16.7.1)} (p. 488), \unsure{d'.} \unsure{qui}, \unsure{devrait} même \unsure{venir avant|venir} \uline{\unsure{lui}} : on \unsure{en prouvera|en première} \unsure{l'unicité|primitifs}, en \unsure{ce particulier|particulier} $\pi$-accessible, \unsure{en disant} 4.13.7 \unsure{vient|venant} après 4.13.6)}
Pour tt $\alpha\in A$
\[
\begin{aligned}
\underline S_\alpha(\pi)&\overset{\mathrm{def}}{=}\underline S_\alpha\cap M_\pi\\
\underline S(\pi)&\overset{\mathrm{def}}{=}\underline S\cap M_\pi=\bigcap\underline S_\alpha(\pi).
\end{aligned}
\]
\deleted{Chacun} \unsure{Chque|Chaque} $\underline S_\alpha(\pi)$ est stable dans $M$
par $\varinjlim$ gdes devant $\pi'_0$ et de cardinal
$\leq\pi$ (puisqu'il en est ainsi de $\underline S_\alpha$ et de
$M_\pi$), [donc il en est de $=$ de leur intersection
\scanpage{377}
\mspage{293}
% archive sheet number: 376

$\underline S(\pi)$. A fortiori $\underline S(\pi)$ est une sous-catégorie
str. pleine et \unsure{Karoubienne|karoubienne} de $M_\pi$.]

\deleted{\unsure{On aura|En outre}, puisque chaque} D'autre part, chaque $\underline S_\alpha$ étant stable \inserted{dans $M$} par
$\varinjlim$ gdes devant $\pi'_0$, il en est de $=$ de
leur intersection $\underline S$. A fortiori, celle-ci
est stable par $\varinjlim$ gdes dev.\textsuperscript{t} $\pi$, donc
le foncteur pl. fidèle canonique
$\operatorname{Ind}_\pi(\underline S(\pi))\unsure{\text{---}|\longrightarrow}M$ se factorise par $S$
en
\[
\operatorname{Ind}_\pi\underline S(\pi)\hookrightarrow\underline S,
\]
il faut prouver que ce foncteur est
\unsure{essent\textsuperscript{ellt}|essent\textsuperscript{llmt}} surjectif. Soit donc $x\in\deleted{\unsure{\operatorname{Ob}}}S=\operatorname{Ob}\underline S$,
tout revient à prouver que
\[
\varinjlim_{i\in\underline S(\pi)_{/x}}x_i\xrightarrow{\sim}x,\tag*{$(*)$}
\]
et que $\underline S(\pi)_{/x}$ est g\textsuperscript{r}de devant $\pi$.

\deleted{Il revient au même de} Notons que, comme $\pi'_0\geq\pi_0$, et \deleted{\unsure{on}}
$\pi^{\pi'_0}=\pi$, on aura $\pi^{\pi_0}=\pi$, donc $\pi$ est
cardinal utile pour $M$, donc \deleted{\unsure{on a}}
$I=M_{\pi/x}$ est g\textsuperscript{r}de dev.\textsuperscript{t} $x$, et on a
\[
\varinjlim_{i\in I=M_{\pi/x}}x_i\xrightarrow{\sim}x.\tag*{$(**)$}
\]
Donc, si l'on admet \inserted{que} $\underline S(\pi)_{/x}$ est \deleted{filtrante}\inserted{g\textsuperscript{r}de dev.\textsuperscript{t} $\pi$},
\scanpage{378}
\mspage{294}
% archive sheet number: 377

le fait que $\unsure{\alpha|(*)}$ soit iso équivaut à dire
que le foncteur \deleted{d'inclusion} canonique
\[
J\overset{\mathrm{def}}{=}\underline S(\pi)_{/x}\hookrightarrow M_{\pi/x}=I
\tag*{\deleted{$\unsure{(*}$}}
\]
est \uline{cofinal}. (cf 4.4.2 p. 79). \deleted{\unsure{[}}Mais
on voit tout de suite que ce foncteur
est un foncteur d'inclusion d'une sous-catégorie \inserted{str.} pleine [, stable par $\varinjlim$ gdes
devant $\pi'_0$, de cardinal $\leq\pi$].

Notons maintenant que la flèche
$(**)'$ similaire à $(*)$, \unsure{où|avec} $\underline S(\pi)$
remplacée par $\underline S_\alpha(\pi)$, est bel et bien
un iso, puisque $x\in\operatorname{Ob}\underline S\subset\operatorname{Ob}\underline S_\alpha$,
et que $\underline S_\alpha$ est $\pi$-accessible. Cela
s'interprète par le fait que le foncteur
d'inclusion
\[
J_\alpha\overset{\mathrm{def}}{=}\underline S_\alpha(\pi)_{/x}\unsure{\hookrightarrow|\lhook\joinrel\relbar}M_{\pi/x}=I
\]
est cofinal. C'est encore une inclusion
\deleted{\unsure{pl.}} d'une sous-catégorie \inserted{$J_\alpha$} strict\textsuperscript{mt} pleine \inserted{de $I$},
stable par $\varinjlim$ g\textsuperscript{r}des devant $\pi'_0$ et de
cardinal $\leq\pi$, et on a
\[
J=\bigcap J_\alpha,
\]
Donc la conclusion voulue résultera du
\scanpage{379}
\mspage{295}
% archive sheet number: 378

% slanted
\margin{\deleted{\illegible{1}} Il faut supposer que $\pi$ \unsure{égal\textsuperscript{t}|égal.} ne soit pas un cardinal limite, pour pouvoir conclure que $I$ est stable par $\varinjlim$ de type $I_{\alpha_\pi}$ et les $\unsure{J_r|J_\alpha}$ stables dans $I$ \unsure{rel\textsuperscript{t} à|pour} ces limites !}
\uline{Lemme} 4.13.10. Soient \inserted{$\pi'_0$ et $\pi>\pi'_0$ des cardinaux infinis,}
$I$ une catégorie équivalente à une petite catégorie,
\deleted{g\textsuperscript{r}de devant $\pi$,} $(J_\alpha)_{\alpha\in A}$ une famille
de sous-catégories \deleted{\unsure{de}} $J_\alpha$ de $I$. On
suppose :

\textcircled{a} $I$ est g\textsuperscript{r}de devant $\pi$.

\margin{Il suffit que $I$ soit g\textsuperscript{r}de devant \deleted{tt} $\pi'<\pi$}

\deleted{a) $\operatorname{card} A\leq\pi$}

\deleted{b) $\forall\alpha\in A$, $J_\alpha$ est une sous-catégorie
\uline{cofinale} de $I$, \uline{g\textsuperscript{r}de devant $\pi$},}

\deleted{\unsure{Soit de plus} \illegible{2}}
\textcircled{b} \deleted{$\forall\alpha\in A$, $J_\alpha$ est stable dans} \inserted{$I$ est} stable par $\varinjlim$ de
cardinal $\leq\pi$, et g\textsuperscript{r}des devant $\pi'_0$.

% margin note written sideways
\margin{Il suffit ici la stabilité par $\varinjlim$ de type $I_{\alpha_\pi}$, cf \unsure{pages|page} \unsure{302|308}}

\textcircled{\unsure{c|b}} $\forall\deleted{\unsure{\alpha}}r\in A$, $J_{\deleted{\unsure{\alpha}}r}$ est stable dans $I$ par
les limites inductives envisagées dans b)
(en particulier $J_{\deleted{\unsure{\alpha}}r}$ est une sous-catégorie
\uline{str. pleine} de $I$).

\textcircled{d} $\forall r\in A$, $J_r$ est une sous-catégorie
\uline{cofinale} de $I$, [\uline{g\textsuperscript{r}de devant} $\pi$].

% slanted
\margin{Cf forme plus générale de ce lemme dans 4.13.10.2 p. \unsure{302|352}}
\textcircled{e} $\operatorname{card}A\leq\pi$.

Soit $J=\bigcap_{r\in A}J_r$. Alors $J$ est une
cofinale dans $I$ \inserted{et} g\textsuperscript{r}de devant $\pi$
(et de plus stable \inserted{dans $I$} par $\varinjlim$ de cardinal $\leq\pi$
et g\textsuperscript{r}des dev.\textsuperscript{t} $\pi'_0$).

En d'autres termes, si $I$ est g\textsuperscript{r}de
dev.\textsuperscript{t} $\pi$, l'intersection \inserted{$J$} d'une famille
$ {J_\alpha}_{\alpha\in A}$ de sous-catégories strict\textsuperscript{mt} pleines
cofinales et g\textsuperscript{r}des dev.\textsuperscript{t} $\pi$, est encore
\scanpage{380}
\mspage{296}
% archive sheet number: 379

cofinale et g\textsuperscript{r}de devant $\pi$, \uline{pourvu que}
$\operatorname{card} A\leq\pi$, et que l'on \deleted{suppose}\inserted{postule} une
propriété de stabilité de $I$ et des
$J_\alpha$ dans $I$, pour les limites inductives
de cardinal $\leq\pi$ et g\textsuperscript{r}des devant un
cardinal convenable $\pi'_0<\pi$ donné.

4.13.10.1. Notons d'ailleurs que si
\deleted{\uline{Démonstration \unsure{du lemme|des lemmes}}}
$f:J\unsure{\text{---}|\longrightarrow}I$ est un foncteur \inserted{cofinal} \uline{plein\unsure{mt}} \uline{fidèle},
avec $I$ g\textsuperscript{r}de devant $\pi$, alors $J$ est
automatiquement g\textsuperscript{r}de devant $\pi$. Cela
résulte p. ex. de 4.9.6.3 $2^\circ$ (p. \unsure{225|235})
par le critère $C_\pi$ (vu que pour $f$ pl fid,
la condition $\unsure{\mathrm E.3_{1\pi}|\mathrm E.3_{17}}$ est tautologiquement
satisfaite). On peut aussi le voir ainsi.
Soit $B$ un ens. avec $\operatorname{card}B\leq\pi$,
il faut prouver que $\deleted{\unsure{B}}\unsure{J|B|I}\unsure{\text{---}|\longrightarrow}J^B$ est cofinal,
or on a

\diagram{\begin{tikzcd}[ampersand replacement=\&]
J \arrow[r,dash,"{\mathrm{cof.}?}"] \arrow[d,"{\mathrm{cof.}}"'] \& J^B \arrow[d,"{\mathrm{cof.}}"]\\
I \arrow[r,dash,"{\mathrm{cof.}}"] \& I^B \mathpunct{,}
\end{tikzcd}}

où trois des quatre flèches du carré
sont cofinales, $J^B\unsure{\text{---}|\longrightarrow}I^B$ l'étant comme
produit de foncteurs cofinaux entre
catégories \uline{filtrantes} (cf 4.9.6.5, page 227)
-- on sait en effet déjà que $J$ est filtrante.
\scanpage{381}
\mspage{297}
% archive sheet number: 380

Donc le composé
\diagram{\begin{tikzcd}[ampersand replacement=\&]
J \arrow[r] \& J^B \arrow[r,"{\mathrm{cof.}}"] \& I^B
\end{tikzcd}}
est cofinal. Mais comme \deleted{il s'agit}\inserted{le deuxième}
\deleted{de} foncteur \inserted{$J^B\unsure{\text{---}|\longrightarrow}I^B$ est} pl. fid., \deleted{la catégorie}
il s'ensuit aussitôt que $J\longrightarrow J^B$
est cofinal (car $\unsure{\xi|\zeta}\backslash J$ pour $\xi\in\operatorname{Ob}J^B$
est le $=$ \inserted{à iso près}, qu'on considère $\xi$ comme
objet de $J^B$, ou de $I^B$).

Ainsi dans l'énoncé de la condition
$d)_{\unsure{r|\alpha}}$ du lemme, il est inutile de
spécifier que $J_r$ est g\textsuperscript{r}de devant $\pi$,
et de $=$ dans la conclusion sur $J$.

Comme $J\hookrightarrow I$ est pl. fidèle,
pour voir que ce foncteur est cofinal,
il suffit de voir qu'il satisfait F1,
i.e. que pour tout $i\in\operatorname{Ob} I$,
existe $j\in J$ avec $i\longrightarrow j$. On
va construire $j$ comme limite d'un
\deleted{suite \unsure{ord.}}\inserted{syst. inductif.} ordin\unsure{al|ale} d'objets \inserted{$i_\alpha$} de $I$,
indexée par des ordinaux $\alpha$ tels que
$\operatorname{card}\alpha<\pi\deleted{\unsure{{}_0|{}'_0}}$.
\deleted{On aura \unsure{donc} \inserted{\unsure{Comme} $\pi'_0<\pi$},} Désignant
par $\alpha_\pi$ le plus petit ordinal tel que
\scanpage{382}
\mspage{298}
% archive sheet number: 381

que $\operatorname{card}\alpha_\pi\deleted{\unsure{>}}=\pi\deleted{\unsure{{}_0|{}'_0}}$,
\deleted{$\operatorname{card}\alpha_0$ est \inserted{\deleted{donc}} le successeur de $\pi'_0$, donc $\leq\pi$.}
L'ens \inserted{ordonné} $I_{\alpha_\pi}$ des
ordinaux $\alpha<\alpha_\pi$ étant g\textsuperscript{r}de devant $\pi'_0$
et de cardinal \deleted{$\unsure{\leq}$} $\pi$, on pourra donc
prendre la limite $\varinjlim_{\alpha\in I_{\alpha_{\unsure{0|\pi}}}}i_\alpha$ dans $I$.
% slanted
\margin{Faux si $\pi$ est un cardinal limite \unsure{!}}

$\circ$ Construction de $(i_\alpha)_{\alpha\in I_{\alpha_0}}$.

\deleted{\unsure{$i_0=i$}} Soit
\[
\begin{aligned}
\alpha_\pi&=\text{plus petit ordinal \inserted{$\alpha$} tel que}\\
&\phantom{=}\operatorname{card}\alpha=\pi
\end{aligned}\tag*{(1)}
\]
donc c'est l'ordinal caractérisé par
\[
\left\{\begin{aligned}
&\operatorname{card}\alpha_\pi=\pi\\
&\alpha<\alpha_\pi\Longrightarrow\operatorname{card}\alpha<\pi,
\end{aligned}\right.\tag*{\unsure{(1)|(2)}}
\]
ou encore par le fait que l'ens ord.
$I_{\alpha_\pi}$ des ordinaux $\alpha<\alpha_\pi$ est \uline{parfaite\textsuperscript{mt}
bien ordonné}. \deleted{\unsure{Donc}} Ainsi, $I_{\alpha_\pi}$ \uline{est
g\textsuperscript{r}d devant $\pi'$ pour tt cardinal} $\pi'<\pi$.
% slanted
\margin{Faux si $\pi$ est cardinal limite \unsure{!}}

Dans les conditions b) \inserted{c)} de 4.13.10, il \unsure{suffit|suffirait}
de supposer $I$ stable par $\varinjlim$ de
type $I_{\alpha_\pi}$, et les $J_{\unsure{r|\alpha}}$ $(r\in A)$
\unsure{stable|stables} dans $I$ par ce type de limites.

\deleted{\unsure{Co}} \uline{Lemme} 4.13.10.1. Soit $\pi$ un cardinal infini,
$A$ un ensemble \inserted{non vide} tel que $\operatorname{card}A\leq\pi$. Alors
\scanpage{383}
\mspage{299}
% archive sheet number: 382

il existe une application
\[
\unsure{\varphi|\phi}:I_{\alpha_\pi}\unsure{\longrightarrow|\text{---}}A
\]
telle que pour tout $r\in A$, $\varphi^{-1}(\{r\})$ soit
une partie cofinale de $I_{\alpha_\pi}$, ou ce
qui \deleted{revient au $=$}\inserted{est plus \deleted{\unsure{fort}}précis},
$\operatorname{card}\varphi^{-1}(\{r\})=\pi$.

En effet il suffit de considérer
un ensemble $A'$ de cardinal $\pi$. (p. ex. $A'=I_{\alpha_\pi}$), et
l'application
\[
\psi=\operatorname{pr}_1:A\times A'\unsure{\text{---}|\longrightarrow}A,
\]
et de se donner une bijection \deleted{de}
\[
\deleted{A\times A}\qquad A\times A'\xrightarrow{\sim}I_{\alpha_\pi},
\]
laquelle existe puisque
\[
\operatorname{card}(A\times A')=\operatorname{card}A\mathbin{.}\operatorname{card}A'=\operatorname{card}A'=\pi\inserted{{}=\operatorname{card}I_{\alpha_\pi}}
\]
puisque $0\neq\operatorname{card}A\leq\operatorname{card}A'=\pi$ et \deleted{\unsure{$\operatorname{card}A'$ infini}}\inserted{$\pi$ infini}.

\deleted{On choisira} On suppose choisie $\varphi$
comme dans 4.13.10.1, et on \deleted{choisi\unsure{ra}}
construira le système inductif
$(i_\alpha)_{\alpha\in\deleted{\unsure{A}}I_{\alpha_\pi}}$ de telle façon qu'on ait
\[
i_\alpha\in J_{\varphi(\alpha)}.\tag*{\deleted{$\unsure{(2)}$} $(*)$}
\]
Posant
\[
j=\varinjlim_{\alpha\in I_{\unsure{\alpha|\alpha_\pi}}}i_\alpha,
\]
\scanpage{384}
\mspage{300}
% archive sheet number: 383

\deleted{\unsure{on aura}}\unsure{pour} tout $r$, comme
\[
I_r\overset{\mathrm{def}}{=}\unsure{\varphi|\phi}^{-1}(r)\subset I_{\alpha_\pi}
\]
est cofinal dans $I_{\alpha_\pi}$, on aura donc
\[
j=\varinjlim_{\alpha\in\deleted{\unsure{\varphi^{-1}(r)}}I_r}i_\alpha,
\]
et \deleted{comme les} en vertu de $(*)$ les
$i_\alpha$ pour $\alpha\in I_r$ sont dans $J_r$, donc
par passage à la limite, utilisant
l'hypothèse c), on trouve que $j\in J_r$,
donc \deleted{$j\in$}
\[
j\in\bigcap_{r\in A}J_r=J,
\]
et comme on choisira
\[
i_0=i,\tag*{$(**)$}
\]
on aura bien trouvé $i\longrightarrow j$ avec $i\in J$.
Donc il suffit de construire le système
inductif $(i_\alpha)$ satisfaisant $(*)$ et $(**)$. 

Si \inserted{le syst\textsuperscript{ème} inductif est construit jusque \unsure{compris|et y compris}} $i_\alpha\in\operatorname{Ob}J_{\varphi(\alpha)}$, \deleted{\unsure{est cons}} on continue
par
\[
i_{\alpha\deleted{\unsure{+1}}}\longrightarrow i_{\alpha+1}\in\operatorname{Ob}J_{\varphi(\alpha+1)},
\]
ce qui est possible puisque $J_{\varphi(\alpha+1)}$ est
\scanpage{385}
\mspage{301}
% archive sheet number: 384

cofinale. Donc il reste à \deleted{définir} construire
$i_\alpha$ pour $\alpha$ ordinal limite. Or
l'ens des $i_\beta$ pour $\beta<\alpha$ est de
cardinal $<\pi$, et comme \inserted{par a)} $I$ est g\textsuperscript{r}de
devant $\pi$ (il suffit qu'il le soit devant
tout cardinal $<\pi$) il existe un
$i\in\operatorname{Ob}I$ qui \unsure{majore} tous les $i_\beta$.
Choisissons des
\[
i_\beta\xrightarrow{\unsure{w|u}_\beta}i,
\]
On n'aura pas a priori comm. des

\diagram{\begin{tikzcd}[ampersand replacement=\&]
\deleted{\unsure{v_\alpha\tau_\beta}}\phantom{i}\\[-1ex]
i_{\beta'}\arrow[r,"{\tau_{\beta,\beta'}}"]\arrow[rr,bend right=17,"{u_{\beta'}}"'] \& i_\beta\arrow[r,"{u_\beta}"] \& i\mathpunct{,}
\end{tikzcd}}

mais comme l'ens des doubles-flèches

\diagram{\begin{tikzcd}[ampersand replacement=\&]
i_{\beta'}\arrow[r,shift left=0.5ex,"{u_\beta\tau_{\beta,\beta'}}"]\arrow[r,dash,shift right=0.5ex,"{u_{\beta'}}"'] \& i
\end{tikzcd}}

est de cardinal $<\pi$, $I$ étant g\textsuperscript{r}de devant $\pi$,
il existe $i\longrightarrow i'$ qui les égalise
tous. Comme $J_{\unsure{\varphi|\phi}(\alpha)}$ est cofinal dans
$I$, il existe $i''\in J_{\unsure{\varphi\alpha|\varphi(\alpha)}}$, et $i'\unsure{\text{---}|\longrightarrow}i''$.
On prend $i_\alpha=i''$, avec les hom composés

\diagram{\begin{tikzcd}[ampersand replacement=\&]
i_\beta \arrow[r,"{u_\beta}"]\arrow[rrr,bend right=12,"{\tau_{\alpha,\beta}}"'] \& \unsure{i_\beta|u_\beta|i}\arrow[r] \& i'\arrow[r,dash] \& i^{\unsure{\prime\prime|\prime}}
\end{tikzcd}}
\scanpage{386}
\mspage{302}
% archive sheet number: 385

comme \deleted{composés} morphismes de transition.
Cela achève la démonstration de 4.13.10,
et par là de 4.13.9, 4.13.8.

\uline{Remarque} 4.13.10.2. Voici ce qu'on a eu
à utiliser vraiment, pour établir la
conclusion de 4.13.10, en termes de $I,\pi$,
$(J_r)_{r\in A}$ (et sans introduire de $\pi'_0$,\unsure{)}\deleted{\unsure{$M$}}

\textcircled{a'} $I$ est g\textsuperscript{r}de devant tout cardinal \inserted{$\pi'$} $<\pi$.

\textcircled{b'} $I$ est stable par $\varinjlim$ de type $I_{\alpha_\pi}$

\textcircled{c'} $\forall r\in A$, $J_r$ est stable dans $I$ par les
$\varinjlim$ de type $I_{\alpha_\pi}$

\textcircled{\textcircled{d'}} $\forall r\in A$, $J_r$ est cofinal dans $I$

\textcircled{e\unsure{'|''}} $\operatorname{card}A\leq\pi$.

(NB. d') et e') ne sont autres que d) et e)
de 4.3.10, tandis que a'), b'), c'), sont
des versions affaiblies de a), b), c)
respectivement.

\uline{4.13.8} (reformulation) Soit pour $r\in A$, $\pi''_r\geq\pi'_r\geq\pi_0$ (= card. utile \inserted{\unsure{pour $M$}} \unsure{min.})
comme spécifié dans 4.13.2 b) (et c), \deleted{\unsure{et}} soit $\pi'_r=\operatorname{Sup}_{r\in A}\pi'_r$, $\pi''\deleted{\unsure{{}_r}}=\operatorname{Sup}_{r\in A}\pi''_r$, de
sorte que pour tout cardinal $\pi$ où $\pi^{\unsure{\pi'|\pi'_r|\pi'_0}}=\pi\geq\pi''$, $\underline S_r$ soit $\pi$-accessible
dans $M$ (et $\pi$ utile pour $\underline S_r$ \uline{et} pour $M$). Soit $c$ un cardinal \uline{non
limite} tel que \deleted{\unsure{$\operatorname{card} A,\operatorname{Sup}(\pi',\pi'')$}}\inserted{$c\geq\operatorname{card}A$ et $c>\pi'$}
\deleted{\unsure{p.ex.} $\unsure{c=\operatorname{card}A}$ si $\operatorname{card} A$ est \unsure{majoré}}
(p.ex. \deleted{limite, ou} le successeur de \deleted{$\unsure{\operatorname{card}(A)}$}\inserted{$\operatorname{Sup}(\operatorname{card}A,\pi')$}.) Prenons $\pi$ tel
que
\[
\pi^{\pi'}=\pi\geq\operatorname{Sup}(\pi'',c).
\]
Alors $\underline S=\bigcap_{r\in A}\underline S_r$ est $\pi$-accessible.

En effet, si $\alpha_c$ est le premier ordinal de cardinal $c$, comme $c$ n'est
pas limite, $I_{\alpha_c}$ est grand devant tt cardinal $c'<c$, en particulier
devant $\pi'$. \deleted{\unsure{$\geq\pi'$}} Or \deleted{\unsure{les $M_\pi$}} pour tt $x\in\operatorname{Ob}M$, $I=M_{\pi/x}$ et les $J_r=\underline S_{r/x}$
sont stables par $\varinjlim$ g\textsuperscript{r}des dev\textsuperscript{t} $\pi'$ et de cardinal $\leq\pi$, \deleted{\unsure{donc} \unsure{les}}
\scanpage{387}\mspage{302 bis}

\deleted{\unsure{devient}} \deleted{\unsure{est}} les $J_r$ stables dans $I$ pour ces limites,
donc il en est en particulier ainsi pour les $\varinjlim$
de type $I_{\alpha_c}$. On peut donc appliquer le résultat
4.13.10.2 pour conclure que si les $J_r$ sont
cofinaux \deleted{\unsure{ds}} $I$ (qui est grd devant $\pi$, $\pi$ étant utile
pour $M$), leur $J=\underline S_{/x}$ l'est également. Cela implique
alors, comme on a vu, que $\underline S$ est $\pi$-accessible
dans $M$\unsure{!|.}

\uline{4.13.10.3.} On peut reformuler la reformulation
précédente, en utilisant la terminologie introduite
plus bas (4.13.23) :

Soit $c$ un cardinal, \deleted{\unsure{non limite}} $\alpha_c$ le plus petit
ordinal $\alpha$ tel que $\operatorname{card}\alpha=c$, $\pi$ un cardinal $\geq c$.

On suppose

$(*)$ \deleted{\unsure{que}} $\pi$ est \inserted{utile} pour $M$, et $M$ est stable \unsure{par|pour} $\varinjlim$ de
type $I_{\alpha_c}$. (\inserted{\deleted{\unsure{utile}} \unsure{admet un card.}}c'est le cas p. ex. si \inserted{$c$ est \uline{non limite} et} $c>\pi_0$, $\pi_0$ étant un
cardinal \deleted{\unsure{utile}} tel que $M$ soit grd $\unsure{\mathrm{dt}|\mathrm{dvt}}$ $\pi_0$, p. ex. un
cardinal \uline{utile} pour $M^{(*)}$.)
% slanted
\margin{$^{(*)}$ Car alors $I_{\alpha_c}$ est grd devant $\pi_0$}

Alors l'ensemble des sous-catégories $\underline S$ de $M$ qui
sont $\alpha_c$-$\pi$-accessibles dans $M$ (i.e. qui sont
$\pi$-accessibles et stables dans $M$ par $\varinjlim$ de type $I_{\alpha_c}$)
est stable par intersections de cardinal $\leq c$. Plus
généralement, si $\unsure{\mathscr L|\underline{\mathscr L}}$ est un ens. de catégories
telles que $1^\circ)$ $I\in\mathscr L\Longrightarrow\operatorname{Card} I\leq\pi$ et $2^\circ)$ $I_{\alpha_c}\in\mathscr L$,
alors l'ens. des sous-catégories de $M$ qui sont
$\mathscr L$-$\pi$-accessibles dans $M$, \unsure{est} stable par \unsure{l'intersection|intersections}
\unsure{de|en} cardinal $\leq c$.
% archive sheet number: 386
\scanpage{388}
% printed form
% letterhead
% A rising diagonal pen stroke cancels each of the two printed columns.
\noindent\begin{minipage}[t]{.48\linewidth}
\deleted{MINISTERE DE L'EDUCATION NATIONALE\\
ACADEMIE D'AIX-MARSEILLE\\
\rule{2cm}{0.3pt}\\
Inspection Académique\\
de Vaucluse\\
\rule{2cm}{0.3pt}\\
Tél. : 90 82 99 76\\
Référence : Secrétariat n°}
\end{minipage}\hfill
\begin{minipage}[t]{.48\linewidth}
\deleted{AVIGNON, le \unsure{.}\\[2em]
L'Inspecteur d'académie,\\
directeur des services départementaux\\
de l'Education Nationale de Vaucluse,\\[2em]
à M}
\end{minipage}
% archive sheet number: 387
